\documentclass[a4paper,10pt]{report}

\usepackage[utf8]{inputenc}
\usepackage[T1]{fontenc}
\usepackage[french]{babel}
\usepackage{lmodern}
\usepackage{tikz}
\usepackage{pgfplots}
\makeatletter
\def\input@path{{../../../../}}
\makeatother
\usepackage[fancyhdr]{teach}
\graphicspath{{../../../../}}

\usetikzlibrary{arrows}
\pgfplotsset{compat=1.18}
\newcolumntype{S}{p{0pt}}

\definecolor{alizarin}{rgb}{0.82, 0.1, 0.26}
\definecolor{meatbrown}{rgb}{0.9, 0.72, 0.23}
\definecolor{olivine}{rgb}{0.6, 0.73, 0.45}

\redefinecolor{thm}{title}{bgcolor}{alizarin}
\redefinecolor{prop}{title}{bgcolor}{meatbrown}
\redefinecolor{defn}{title}{bgcolor}{olivine}

\redefinetitle{ex}{Application}
\redefinetitle{rem}{Remarque}
\redefinetitle{memento}{Règle}

\newlength{\blociconewidth}
\setlength{\blociconewidth}{8mm}
\newcommand{\teachblocicon}[1]{\raisebox{-2mm}{\includegraphics[width=6mm]{#1}}}

\makeatletter
\renewenvironment{memento}[1][]
{%
\ifx\hfuzz#1\hfuzz
\renewcommand\teach@title{\teach@title@memento}
\else
\renewcommand\teach@title{\teach@title@memento~(#1)}
\fi
\begin{lrbox}{\teach@box}
\begin{minipage}{\dimexpr\linewidth-\blociconewidth-\arrayhspace-\arrayvertrulewidth-2\tabcolsep}
\color{\teach@thm@body@txtcolor}
}
{
\end{minipage}
\end{lrbox}
\vspace{1em}
\begin{tabular}{@{}p{\blociconewidth}@{\color{\teach@memento@title@txtcolor}\hspace*{1mm}\vrule width \arrayvertrulewidth\hspace*{\arrayhspace}}p{\dimexpr\linewidth-\blociconewidth-\arrayhspace-\arrayvertrulewidth-2\tabcolsep}@{}}
\teachblocicon{memento.png} & \raisebox{1mm}{\fontfamily{lmss}\bfseries\selectfont\textcolor{\teach@memento@title@txtcolor}{\teach@title}}\\[-1pt]
 & \usebox{\teach@box}
\end{tabular}
\vspace{1em}
}

\renewenvironment{attention}[1][]
{%
\ifx\hfuzz#1\hfuzz
\renewcommand\teach@title{\teach@title@attention}
\else
\renewcommand\teach@title{\teach@title@attention~(#1)}
\fi
\begin{lrbox}{\teach@box}
\begin{minipage}{\dimexpr\linewidth-\blociconewidth-\arrayhspace-\arrayvertrulewidth-2\tabcolsep}
\color{\teach@thm@body@txtcolor}
}
{
\end{minipage}
\end{lrbox}
\vspace{1em}
\begin{tabular}{@{}p{\blociconewidth}@{\color{\teach@attention@title@txtcolor}\hspace*{1mm}\vrule width \arrayvertrulewidth\hspace*{\arrayhspace}}p{\dimexpr\linewidth-\blociconewidth-\arrayhspace-\arrayvertrulewidth-2\tabcolsep}@{}}
\teachblocicon{attention_img.png} & \raisebox{1mm}{\fontfamily{lmss}\bfseries\selectfont\textcolor{\teach@attention@title@txtcolor}{\teach@title}}\\[-1pt]
 & \usebox{\teach@box}
\end{tabular}
\vspace{1em}
}
\makeatother

\newcommand{\grandscarreaux}[2]{%
  \begin{tikzpicture}
    \draw[step=0.8cm,gray!45,very thin] (0,0) grid (#1,#2);
  \end{tikzpicture}%
}

\begin{document}

\newpage
\setcounter{chapter}{0}
\chapter{Nombres entiers}

\section{Décomposition, nom des chiffres}

\begin{memento}
0, 1, 2, 3, 4, 5, 6, 7, 8 et 9 sont des chiffres qui permettent d'écrire tous les nombres entiers, de même que les lettres de A à Z permettent d'écrire tous les mots.
\end{memento}

\begin{ex}
\noindent\grandscarreaux{\linewidth}{2.4cm}
\end{ex}

\begin{memento}
Pour pouvoir lire les grands nombres entiers facilement, on regroupe les chiffres par  \dotfill
\end{memento}

\begin{ex}

\medskip On s'intéresse au nombre $1048774912$.
\begin{enumerate}
\item Réécris ce nombre en appliquant la règle précédente.
\item Écris ce nombre en toutes lettres.
\item Décompose ce nombre
\item Donne le nom des chiffres 8 et 9.
\item Quel est le nombre des millions de ce nombre ?
\end{enumerate} 
\noindent\grandscarreaux{\linewidth}{6.4cm}
\end{ex}

\section{Repérage sur une demi-droite graduée}

\begin{defn}
Une demi-droite graduée est une demi-droite sur laquelle on a reporté une unité de longueur régulièrement (souvent le centimètre) à partir de son origine.
\end{defn}

\begin{prop}
Sur une demi-droite graduée, un point est repéré par un nombre appelé son \dotfill

\medskip \noindent L'origine est repéré par le nombre \dots \dots .
\end{prop}

\begin{ex}
\medskip

\begin{center}
\begin{tikzpicture}[line cap=round,line join=round,>=triangle 45,x=0.045cm,y=1.0cm]
\begin{axis}[x=0.045cm,y=1.0cm,axis lines=middle,ymajorgrids=true,xmajorgrids=true,
xmin=-1.0,xmax=345.0,ymin=-0.1,ymax=0.1,xtick={0.0,20.0,...,340.0},ytick={-0.0,20.0,...,0.0},]
\clip(-1.,-1) rectangle (345.,1);
\begin{scriptsize}
\draw [color=blue] (60.,0.)-- ++(-3.5pt,-3.5pt) -- ++(7.0pt,7.0pt) ++(-7.0pt,0) -- ++(7.0pt,-7.0pt);
\draw[color=blue] (62.67583814733487,1) node {$A$};
\draw [color=blue] (110.,0.)-- ++(-3.5pt,-3.5pt) -- ++(7.0pt,7.0pt) ++(-7.0pt,0) -- ++(7.0pt,-7.0pt);
\draw[color=blue] (112.6420565791818,1) node {$B$};
\draw [color=blue] (310.,0.)-- ++(-3.5pt,-3.5pt) -- ++(7.0pt,7.0pt) ++(-7.0pt,0) -- ++(7.0pt,-7.0pt);
\draw[color=blue] (312.76014081402155,1) node {$C$};
\end{scriptsize}
\end{axis}
\end{tikzpicture}
\end{center}

\noindent\grandscarreaux{\linewidth}{2.8cm}
\end{ex}

\section{Comparaison et rangement}

\begin{defn}
Comparer deux nombres, c'est \dotfill

\medskip\noindent\dotfill
\end{defn}

\begin{defn}
Ranger des nombres dans l'ordre croissant signifie \dotfill

\medskip\noindent\dotfill
\end{defn}

\begin{defn}
Ranger des nombres dans l'ordre décroissant signifie \dotfill

\medskip\noindent\dotfill
\end{defn}

\begin{ex}
Range les nombres 25 342 ; 253 420 ; 25 243 ; 235 420 ; 25 324 dans l'ordre croissant.

\medskip\noindent\grandscarreaux{\linewidth}{2.8cm}
\end{ex}

\section{Addition}

\begin{defn}
Les nombres que l'on additionne s'appellent les \dotfill

\medskip \noindent Le résultat d'une addition s'appelle la \dotfill
\end{defn}

\begin{ex}
Pose et calcule 1 856 + 525.

\medskip\noindent\grandscarreaux{\linewidth}{2.6cm}
\end{ex}

\begin{prop}
Dans une addition, on a le droit de :
\begin{enumerate}
\item regrouper les termes ;
\item changer des termes de place.
\end{enumerate} 
\end{prop}

\begin{ex}
Calcule astucieusement 46 + 37 + 54 + 63.

\medskip\noindent\grandscarreaux{\linewidth}{2cm}
\end{ex}

\section{Soustraction}

\begin{defn}
Les nombres que l'on soustrait s'appellent les \dotfill

\medskip \noindent Le résultat d'une soustraction s'appelle la \dotfill
\end{defn}

\begin{ex}
Pose et calcule 233-67.

\medskip\noindent\grandscarreaux{\linewidth}{2.6cm}
\end{ex}

\begin{attention}
On ne peut pas changer les termes de place dans une soustraction.
\end{attention}

\end{document}
